\chapter{eBPF}

\section{Introduction and overview}

TODO

\section{Instruction Set}

eBPF programs are written in the eBPF \textit{instruction set architecture}
(ISA), often referred to simply as BPF ISA. The ISA defines the set of
instructions that an eBPF program can use, and are expected to be available from
an eBPF runtime. A compiler may be used to generate the BPF instructions from an
higher-level language such as C or Rust to make the development of eBPF program
easier for developers.

Standardizing the BPF ISA is significant to allow compiler developers, tools
developers, eBPF runtime developers and ASIC designers to agree to a common
reference. The formal IETF standardization effort was started from earlier
drafts in October 2023 with RFC 9669; the RFC was finally standardized one year
later and updated multiple times since the first version. All the information
presented in this section are referring to the RFC as last updated in 20 May
2026.

Controversially, the RFC does not standardize the platform-specific application
binary interface (ABI) because a common consensus was not reached and it was
preferred to be left out of the standard. This includes various details such as:
which registers are saved across a call, which registers contain the frame
pointer, how large the stack is, and other specifics. Further standardization is
left for "future work"; however, the Linux BPF ABI is well explained in the
kernel documentation and any BPF runtime that wishes to be compatible with the
Linux implementation must follow it.

The figure below shows C code on the left and an human readable form of BPF
bytecode on the right compiled with clang 23.1.0:

\noindent
\vspace{0.5em}

\begin{minipage}[t]{0.47\textwidth}
\begin{verbatim}
int sum_foo(int a, int b) {
    return foo(a + b);
}
\end{verbatim}
\end{minipage}%
\hfill
\vrule % This creates the vertical dividing bar
\hfill
\begin{minipage}[t]{0.47\textwidth}
\begin{verbatim}
sum_foo(int, int):
        *(u32 *)(r10 - 4) = w1
        *(u32 *)(r10 - 8) = w2
        w1 = *(u32 *)(r10 - 4)
        w2 = *(u32 *)(r10 - 8)
        w1 += w2
        call foo(int)
        exit
\end{verbatim}
\end{minipage}

\noindent
\vspace{0.5em}

Each line on the right corresponds to a BPF assembly instruction. The term
\verb|w1|, \verb|w2| represent the first word of the register number 1 and 2
respectively, and \verb|r10| is the full register number 10.

\subsection*{Instructions and encoding}

The standard specifies two instruction encodings: one using 64 bits and another
one using 128 bits. Instructions are assigned to a named conformance group,
which is a set of instructions or other requirements identified with a short
ASCII label, such as "base32". An implementation of the BPF ISA must support the
"base32" conformance group, but may or may not support other groups.

\begin{figure}[htbp]
    \centering
    \begin{bytefield}[bitwidth=1.1em]{32}
        % Header bit numbers (optional)
        \bitheader{0, 7, 15, 31} \\

        % First row
        \bitbox{8}{opcode} & \bitbox{8}{regs} & \bitbox{16}{offset} \\

        % Second row
        \bitbox{32}{imm}
    \end{bytefield}
    \caption{64 bit instruction encoding.}
    \label{fig:instruction-encoding}
\end{figure}

\begin{figure}[htbp]
    \centering
    \begin{bytefield}[bitwidth=1.1em]{32}
        % Optional: Bit header numbers
        \bitheader{0, 7, 15, 31} \\

        % Row 1: opcode, regs, offset
        \bitbox{8}{opcode} & \bitbox{8}{regs} & \bitbox{16}{offset} \\

        % Row 2: imm
        \bitbox{32}{imm} \\

        % Row 3: reserved
        \bitbox{32}{reserved} \\

        % Row 4: next_imm (remember to escape the underscore!)
        \bitbox{32}{next\_imm}
    \end{bytefield}
    \caption{Extended 128-bit instruction encoding.}
    \label{fig:extended-instruction-encoding}
\end{figure}

Figure \ref{fig:instruction-encoding} shows the encoding of 64 bit instructions.
The various field are the following:

\begin{itemize}

  \item \verb|opcode|: the operation to perform, encoded as a \textit{class}
    part and a \textit{specific} part which identifies the instruction on the
    class.

  \item \verb|regs|: source and destination register numbers. Some instructions
    might reuse this fields for other purposes.

  \item \verb|offset|: signed integer offset used with pointer arithmetic, or
    for other purposes.

  \item \verb|imm|: signed integer immediate value.

\end{itemize}

Figure \ref{fig:extended-instruction-encoding} shows the encoding for extended
128 bit instructions. Compared to 64 bit instructions, it has a \verb|reserved|
32 bit field which is set to 0, and a \verb|next_imm| 32 bit field for a second
signed integer immediate value.

The BTF ISA supports atomic operations for load and store instructions, which
ensure that the operation cannot be interrupted or corrupted by other accesses
to the same memory region by another BPF program or other actors. Other than
atomic \verb|ADD|, \verb|OR|, \verb|AND| and \verb|XOR|, BPF provide \verb|XCHG|
for atomic exchange, and \verb|CMPXCHG| for atomic compare and exchange.

A complete list of the encoding of each instruction is found in the RFC.

\subsection*{Helper Functions}

The RFC introduces the concept of \textit{helper functions} which provide
functionalities implemented by the underlying platform. Notably, helper
functions are not standardized by RFC9669, but they are documented in the Linux
kernel, which will be discussed in a later chapter. To call an helper function,
the jump instruction must be executed with its source register as $0x0$ or
$0x2$, with a platform-specific \verb|ID| identifying a specific function.

\subsection*{Maps}

RFC9669 defines "maps" as shared memory regions accessible by BPF programs on
some platforms. The particular semantics of a map are not defined by the
standard. Linux supports many kinds of maps which implement various data
structures such as hash maps, arrays, ring buffers etc which provide shared
storage across userspace and kernelspace, and among multiple executions of the
same eBPF program or multiple eBPF programs. Linux maps will be discussed later.

The standard also defines "platform variables" which are memory regions
identified by integer ids exposed by the runtime. They are different than maps
as they are used directly by the program instead of going through helper
functions, conceptually platform variables can be thought of as global
variables.

\subsection*{Compiler Support}

Clang 11+ and GCC 12.1+ support compiling C to BPF bytecode. Compilers for other
languages also support the BPF ISA, such as the rust compiler. Atomic operations
which map to atomic instructions can be specified with intrinsics.

\section{The verifier}

After loading and eBPF program, the BPF code in \textit{just in time} (JIT)
compiled to native machine instructions to be executed as kernel code in a later
moment. This would give eBPF programs the ability to run almost arbitrary code
in kernelspace, with important security implications on confidentiality,
integrity and availability. Without a proper security model, an eBPF program
could potentially read or write arbitrary kernel memory (confidentiality and
integrity issue) or invoke other kernel functions and alter the execution of the
operating system to halt it (availability issue). For these reasons, the eBPF
security model determines strict guardrails and limitations to what a program is
allowed to do. The model discriminates a program as \textit{safe} if it meets
all security requirements, or as \textit{unsafe} otherwise; the user can decide
to disallow the loading of any unsafe programs, which is the default policy on
most Linux distributions. This model is fundamental to the success and wide
adoption of eBPF in production and security-critical environments. Other
security considerations and policies will be discussed in a later section, such
as user permissions and chain of trust. Notably, the safety requirements are not
standardized and may change across kernel versions. The reason is that the eBPF
ecosystem is rapidly evolving, hence putting effort in standardization is not
practical or wanted in the current state.

The following is a non-exhaustive list of security requirements for safe eBPF
programs:

\begin{itemize}

  \item Programs must always terminate: no infinite loops or infinite recursion
    is allowed (availability issue).

  \item Programs cannot to read or write arbitrary memory (confidentiality,
    integrity issue).

  \item Network programs are not allowed to access memory outside of packet
    bounds (confidentiality, integrity issue).

  \item Programs are not allowed to deadlock (availability issue).

  \item Programs are not allowed to read uninitialized memory (confidentiality
    issue).

\end{itemize}

These requirements are checked at program load time instead of at runtime, to
trade runtime performance penalties with load-time overhead. If a requirement
cannot be fully checked at load-time, the eBPF program is rejected or a method
to perform the checks dynamically is provided, although only for a limited set
of functionalities (dynamic pointers). The program that performs these checks is
known as the eBPF \textit{verifier}: given a BPF program, the output of the
verifier will determine if the program is considered safe or not. Therefore, the
verifier is a critical part of the eBPF architecture and its implementation has
been widely scrutinized. Given the complexity of checking some of the
requirements at load time, the verifier implementation is highly complex. Just
the file \textit{\$LINUX/kernel/bpf/verifier.c}, which implements most of the
logic, is about $20000$ lines of code long, although solutions to simplify and
modularize it are being discussed by its developers.


\section{Program Types}

TODO

\section{Helper functions}

TODO

\section{Maps}

TODO

\section{BTF and CO-RE}

TODO

\section{Security Considerations}

TODO

\section{Use Cases and modern ecosystem}

TODO
